% Estrutura e diagramação adaptadas do template fornecido pelo pesquisador.
% Referências autor-data com o mesmo backend BibTeX do exemplo de Ana.
\documentclass[12pt,a4paper,twoside]{article}
\usepackage[left=3cm,top=3cm,right=2cm,bottom=2cm]{geometry}
\usepackage[utf8]{inputenc}
\usepackage[T1]{fontenc}
\usepackage[brazilian]{babel}
\usepackage{amsmath,amssymb,bm}
\usepackage{graphicx,booktabs,array,longtable}
\usepackage{fancyhdr,titletoc,titlesec,setspace,indentfirst}
\usepackage{caption,float,placeins}
\usepackage[colorlinks=true,allcolors=black,bookmarksnumbered=true]{hyperref}
\usepackage{bookmark}
\usepackage[alf]{abntex2cite}
\numberwithin{equation}{subsection}
\newcommand{\titulo}{Contribuição preditiva de representações de letras para a popularidade no Spotify}
\newcommand{\autor}{Felipe Macedo Cordeiro}
\newcommand{\orientador}{Leandro Tavares Correia}
\newcommand{\fonte}[1]{\par\vspace{2pt}{\footnotesize Fonte: #1\par}}
\newcommand{\nota}[1]{\par{\footnotesize Nota: #1\par}}
\newcommand{\figura}[4]{\begin{figure}[htbp]\centering\includegraphics[width=\linewidth]{images/#1.png}\caption{#2}\label{#3}\fonte{#4}\end{figure}}
\newcolumntype{P}[1]{>{\raggedright\arraybackslash}p{#1}}
\pagestyle{fancy}\fancyhf{}
\setlength{\headheight}{16pt}
\fancyhead[RO,LE]{\thepage}
\fancyhead[RE,LO]{\nouppercase{\emph{\leftmark}}}
\renewcommand{\sectionmark}[1]{\markboth{#1}{}}
\titlecontents{section}[0cm]{}{\bfseries\thecontentslabel\ }{}{\titlerule*[.75pc]{.}\contentspage}
\titlecontents{subsection}[.75cm]{}{\thecontentslabel\ }{}{\titlerule*[.75pc]{.}\contentspage}
\setcounter{secnumdepth}{3}\setcounter{tocdepth}{2}
\captionsetup{font=small,labelfont=normalfont,justification=centering}
\setlength{\parindent}{1.5cm}\setlength{\parskip}{0.2cm}
\setlength{\intextsep}{0.5cm}
\titlespacing*{\section}{0cm}{0cm}{0.5cm}
\titlespacing*{\subsection}{0cm}{0.5cm}{0.35cm}
\titlespacing*{\subsubsection}{0cm}{0.4cm}{0.25cm}
\raggedbottom
\hypersetup{pdftitle={Contribuição preditiva de representações de letras para a popularidade no Spotify},pdfauthor={Felipe Macedo Cordeiro}}
\begin{document}
\hypersetup{pageanchor=false}
\begin{titlepage}
\centering
\includegraphics[scale=.8]{imagens/unb.png}\par
\vspace{.4cm}
\textbf{Universidade de Brasília\\Departamento de Estatística}\par
\vspace{4.2cm}
\textbf{\titulo}\par
\vspace{2cm}
\textbf{\autor}\par
\vspace{1.5cm}
\begin{flushright}
\begin{minipage}{7.5cm}
Relatório apresentado ao Departamento de Estatística da Universidade de Brasília como parte dos requisitos necessários para obtenção do grau de Bacharel em Estatística.
\end{minipage}
\end{flushright}
\vfill
\textbf{Brasília\\2026}
\end{titlepage}

\thispagestyle{empty}
\begin{center}
\textbf{\autor}\par
\vspace{4cm}
\textbf{\titulo}\par
\vspace{2cm}
Orientador: \orientador
\end{center}
\vspace{2cm}
\begin{flushright}
\begin{minipage}{7.5cm}
Relatório apresentado ao Departamento de Estatística da Universidade de Brasília como parte dos requisitos necessários para obtenção do grau de Bacharel em Estatística.
\end{minipage}
\end{flushright}
\vfill
\begin{center}\textbf{Brasília\\2026}\end{center}

\clearpage\pagenumbering{roman}\hypersetup{pageanchor=true}
\onehalfspacing
\section*{Resumo}\markboth{Resumo}{}
Este trabalho avalia a contribuição preditiva incremental de representações de letras para estimar a popularidade observada no Spotify, em conjunto com atributos acústicos e metadados. O ponto de partida é um catálogo público com 1.159.764 faixas e 82 gêneros. Desse catálogo foi selecionada uma amostra probabilística estratificada de 200.000 faixas, incluindo integralmente as 499 com popularidade de pelo menos 80. A recuperação local de letras no LRCLIB resultou em 77.535 faixas aceitas. A comparação descritiva revelou disponibilidade textual seletiva: a popularidade média do catálogo é 18,38, enquanto a média das aceitas é 23,98. A modelagem foi, portanto, delimitada ao domínio analítico observado, sem pressupor representatividade do catálogo. Compararam-se quatro configurações de regressão Ridge: atributos acústicos e metadados; esses atributos acrescidos de 21 descritores lexicais e estruturais; acrescidos de embeddings multilingues; e acrescidos de ambas as representações textuais. A avaliação utilizou validação cruzada agrupada aninhada, com cinco partições externas e três internas, preservando 21.893 componentes operacionais de artista, composição, versão e texto. O erro médio absoluto passou de 8,3525 no Baseline para 8,1151 no Combinado, redução de 2,84\%. O acréscimo dos descritores aos embeddings produziu ganho menor, de 0,0134 ponto, inconclusivo na análise auxiliar de incerteza. As 128 faixas aceitas com popularidade de pelo menos 80 foram todas subestimadas. Os resultados indicam contribuição textual modesta e não uniforme na avaliação realizada; não demonstram causalidade, previsão prospectiva de sucesso ou generalização automática ao catálogo Spotify.

\vspace{.6cm}
{\raggedright\noindent\textbf{Palavras-chave:} popularidade musical; letras; amostragem estratificada; regressão Ridge; embeddings; validação cruzada agrupada.\par}

\clearpage\listoftables
\clearpage\listoffigures
\clearpage\tableofcontents
\clearpage\pagenumbering{arabic}
\section{Introdução}\label{sec:introducao}
A popularidade de uma música pode ser examinada por medidas distintas, como posição em paradas, número de execuções e escores calculados por plataformas digitais. Essas respostas não são intercambiáveis: dependem do serviço, do período e da unidade observada. No Spotify, a popularidade é um escore de zero a cem, relacionado principalmente às execuções e à sua recência; diferentes versões de uma música podem receber valores distintos \cite{spotifydocs}. Neste estudo, interessa o escore registrado no arquivo utilizado, e não a previsão de popularidade futura ou de sucesso no lançamento.

A análise computacional da música costuma utilizar características do áudio, mas a letra oferece outra fonte de informação. Extensão, diversidade lexical e repetição podem ser descritas por medidas explícitas, enquanto representações vetoriais produzidas por modelos de linguagem sintetizam outras relações entre textos. A questão empírica é se essas representações melhoram a estimação da popularidade quando atributos acústicos e metadados já estão disponíveis.

Esse interesse tem antecedentes na recuperação de informação musical. \citeonline{dhanaraj2005} investigaram a detecção de sucessos com características acústicas e textuais. \citeonline{singhi2014} examinaram características de rima, sílabas e métrica para classificar presença em uma parada musical. Esses problemas de classificação diferem da regressão de um escore contínuo e não permitem transferir diretamente medidas de desempenho para o presente estudo. A análise de mudanças em propriedades de letras realizada por \citeonline{parada2024lyrics}, acompanhada de sua correção \cite{parada2024correction}, também motiva a investigação de descritores, sem estabelecer que eles expliquem ou determinem popularidade.

Uma dificuldade anterior ao ajuste dos modelos é a construção da base analítica. As letras não estão presentes no arquivo Spotify e sua recuperação depende de outra fonte e de regras de correspondência. Selecionar apenas músicas cuja letra foi encontrada pode modificar a distribuição de popularidade, gênero e características acústicas. Assim, a pergunta sobre ganho preditivo precisa ser acompanhada de outra: a quais faixas os resultados podem se referir?

O percurso deste trabalho parte da caracterização de 1.159.764 faixas do catálogo acessível, passa pelo sorteio estratificado de 200.000 faixas e pela recuperação de letras em um snapshot local do LRCLIB, e chega a um domínio observado de 77.535 faixas com letras aceitas. O catálogo acessível não é tratado como censo do Spotify; tampouco o subconjunto com letras é presumido representativo de todas as faixas do catálogo. Essas distinções orientam tanto a descrição dos dados quanto a interpretação dos modelos.

\subsection{Objetivos}
O objetivo geral é avaliar a associação e a contribuição preditiva incremental de representações de letras para estimar a popularidade Spotify observada, em conjunto com atributos acústicos e metadados. Especificamente, busca-se caracterizar as fontes e o efeito da disponibilidade textual sobre a composição da base; documentar a amostragem e a integração; comparar descritores lexicais e estruturais com embeddings multilingues; e avaliar os ganhos nas mesmas observações e divisões de treinamento e teste, incluindo cauda de popularidade e subgrupos linguísticos.

A comparação utiliza a mesma família de regressão, Ridge, em quatro configurações: Baseline, Tradicional, Profundo e Combinado. A denominação Profundo se refere à origem dos embeddings; não indica uma rede neural treinada neste estudo para predizer popularidade. O erro médio absoluto é a métrica principal. O interesse é a contribuição adicional das representações dentro desse procedimento, sem atribuição causal.

Além desta introdução, a Seção~\ref{sec:teoria} apresenta os fundamentos estatísticos e computacionais. A Seção~\ref{sec:metodologia} descreve as fontes, o catálogo inicial, a amostragem, a integração e a base dos modelos. A Seção~\ref{sec:resultados} discute os resultados preditivos e seus diagnósticos. A Seção~\ref{sec:conclusao} reúne as conclusões e os limites de generalização.

\clearpage\section{Referencial teórico}\label{sec:teoria}
\subsection{Amostragem estratificada e estimação}
Considere uma população finita $U$, de tamanho $N$, particionada em $H$ estratos disjuntos $U_h$, com tamanhos $N_h$. Em cada estrato, uma amostra aleatória simples sem reposição $s_h$ contém $n_h$ unidades. A estratificação permite controlar a composição do sorteio por características conhecidas da população; não exige frações de seleção iguais entre estratos. Esses fundamentos são tratados por \citeonline{cochran1977} e \citeonline{lohr2022}.

Para a unidade $i$ pertencente ao estrato $h$, a probabilidade de inclusão e o peso de desenho são
\begin{equation}\label{eq:inclusao}
 \pi_i=\frac{n_h}{N_h},\qquad w_i=\frac{1}{\pi_i}=\frac{N_h}{n_h}.
\end{equation}
Um estrato incluído integralmente tem $n_h=N_h$ e peso um. Para uma variável $z_i$, a média populacional pode ser estimada por
\begin{equation}\label{eq:mediaestratificada}
 \widehat{\overline Z}_{\mathrm{est}}
 =\sum_{h=1}^{H}\frac{N_h}{N}\overline z_h
 =\frac{1}{N}\sum_{i\in s}w_i z_i,
 \qquad\overline z_h=\frac{1}{n_h}\sum_{i\in s_h}z_i.
\end{equation}
Nessa alocação fixa, $\sum_{i\in s}w_i=N$. Sob sorteios independentes entre estratos, a variância estimada da média é
\begin{equation}\label{eq:varianciaestratificada}
 \widehat{\operatorname{Var}}_d(\widehat{\overline Z}_{\mathrm{est}})
 =\sum_{h=1}^{H}\left(\frac{N_h}{N}\right)^2
 \left(1-\frac{n_h}{N_h}\right)\frac{s_{z,h}^2}{n_h},
\end{equation}
em que $s_{z,h}^2$ é a variância amostral no estrato; estratos integrais contribuem com zero. A correção para população finita é relevante quando a fração amostral não é desprezível \cite{cochran1977}. Essas expressões descrevem o sorteio inicial, não o mecanismo posterior de recuperação de letras.

Se $A_i$ indica letra aceita, a média ponderada entre aceitas assume a forma de razão
\begin{equation}\label{eq:mediaaceitas}
 \widehat{\overline Z}_A=\frac{\sum_{i\in s}w_i A_i z_i}{\sum_{i\in s}w_i A_i}.
\end{equation}
Ela descreve a composição do domínio aceito segundo os pesos originais, mas não estima automaticamente a média de $z$ em todo $U$. O evento $A_i=1$ não é uma segunda seleção probabilística com probabilidades conhecidas. Disponibilidade, ambiguidade e qualidade de correspondência devem ser distinguidas de inclusão amostral, em consonância com a discussão de cobertura e não resposta de \citeonline{lohr2022}.

Para a razão da Equação~\ref{eq:mediaaceitas}, quando a variável está observada nos casos considerados, a linearização usa $u_i=A_i(z_i-\widehat{\overline Z}_A)$, incluindo zeros para as não aceitas. Se $\widehat N_A=\sum_{i\in s}w_iA_i$, a variância aproximada é
\begin{equation}\label{eq:varianciarazao}
 \widehat{\operatorname{Var}}_d(\widehat{\overline Z}_A)=\sum_h v_h,
 \qquad v_h=\frac{N_h^2(1-n_h/N_h)s_{u,h}^2}{n_h\widehat N_A^2}.
\end{equation}
Essa é uma aproximação de razão sob o desenho, distinta de tratar as aceitas como um novo sorteio aleatório simples \cite{statcan2010}. Nas comparações marginais, a média do catálogo é fixa, e a implementação utiliza graus de liberdade de Satterthwaite,
\begin{equation}\label{eq:dfrazao}
 \nu=\frac{(\sum_hv_h)^2}{\sum_{h:n_h<N_h}v_h^2/(n_h-1)},
\end{equation}
quando a variância é estimável e positiva. Células sem suporte para essa estimação não recebem conclusão automática de equivalência.

Compatibilidade marginal não se demonstra por deixar de rejeitar uma igualdade. Para uma diferença populacional $\theta$ e margem substantiva $\varepsilon>0$, a formulação de equivalência é
\begin{equation}\label{eq:equivalencia}
 H_0:\theta\le-\varepsilon\ \text{ou}\ \theta\ge\varepsilon,
 \qquad H_1:-\varepsilon<\theta<\varepsilon.
\end{equation}
No procedimento de dois testes unilaterais (TOST) com aproximação por $t$ de Student e nível $\eta$ em cada teste, a equivalência exige que
\begin{equation}\label{eq:tost}
 [\widehat\theta-t_{1-\eta,\nu}\widehat{\operatorname{SE}},
  \widehat\theta+t_{1-\eta,\nu}\widehat{\operatorname{SE}}]
 \subset(-\varepsilon,\varepsilon).
\end{equation}
Aqui $t_{q,\nu}$ é o quantil $q$ com $\nu$ graus de liberdade. Para $\eta=0,05$, usa-se intervalo marginal de 90\%. A lógica de equivalência e a necessidade de justificar margens são discutidas por \citeonline{wellek2010}. Equivalência, diferença apoiada e inconclusão são resultados distintos; concordância em algumas marginais não certifica representatividade multivariada ou identidade correta de letras.

\subsection{Regressão linear e regularização Ridge}
Seja $y_i$ a popularidade observada e $\bm x_i$ o vetor de $p$ características de uma faixa. A regressão linear aproxima a resposta por $\beta_0+\bm x_i^\top\bm\beta$. Com várias características correlacionadas, a estimação por mínimos quadrados pode produzir coeficientes instáveis. A regressão Ridge acrescenta uma penalização quadrática, introduzindo uma troca entre viés e variância \cite{hastie2009,james2023}.

Para o conjunto de treinamento $T$, os coeficientes são estimados por
\begin{equation}\label{eq:ridge}
 (\widehat\beta_0,\widehat{\bm\beta})=
 \underset{\beta_0,\bm\beta}{\operatorname{argmin}}
 \left\{\sum_{i\in T}(y_i-\beta_0-\bm x_i^\top\bm\beta)^2
 +\alpha\sum_{j=1}^{p}\beta_j^2\right\},\qquad\alpha>0.
\end{equation}
O intercepto não é penalizado. O parâmetro $\alpha$ controla o encolhimento dos coeficientes; valores maiores impõem penalização mais intensa. A proposta clássica de Ridge é de \citeonline{hoerl1970}. A comparação aqui mantém essa família nos quatro conjuntos de variáveis para examinar a informação incremental das representações.

Como a penalização depende da escala, cada variável numérica é padronizada no respectivo treino:
\begin{equation}\label{eq:padronizacao}
 x_{ij}^{*}=\frac{x_{ij}-\overline x_{j,T}}{s_{j,T}},\qquad
 s_{j,T}^2=\frac{1}{|T|}\sum_{i\in T}(x_{ij}-\overline x_{j,T})^2.
\end{equation}
A implementação usa desvio com divisor $|T|$. Uma coluna constante é centrada para zero. Médias, desvios e demais transformações são estimados apenas em $T$ e aplicados à avaliação, evitando que informação do teste participe do ajuste \cite{james2023}. Ridge otimiza quadrados penalizados; a escolha de sua penalização pode, contudo, utilizar outra métrica de validação, como o MAE.

\subsection{Medidas de erro e contribuição incremental}
Para $n$ observações de avaliação, denote por $\widehat y_i^{(m)}$ a previsão da configuração $m$, produzida sem usar o componente de $i$ no treino. Defina o erro de previsão $e_i^{(m)}=y_i-\widehat y_i^{(m)}$. O erro médio absoluto (MAE, \textit{mean absolute error}) e a raiz do erro quadrático médio (RMSE, \textit{root mean squared error}) são
\begin{align}
 \operatorname{MAE}_m&=\frac{1}{n}\sum_{i=1}^{n}|e_i^{(m)}|,\label{eq:mae}\\
 \operatorname{RMSE}_m&=\sqrt{\frac{1}{n}\sum_{i=1}^{n}(e_i^{(m)})^2}.\label{eq:rmse}
\end{align}
Essas definições são apresentadas por \citeonline{hyndman2021}. Ambas têm a unidade da resposta; o RMSE é mais sensível a desvios grandes devido ao quadrado. O MAE foi adotado como métrica primária de seleção e comparação.

Também é calculado o coeficiente de determinação nas previsões externas reunidas:
\begin{equation}\label{eq:r2}
 R_m^2=1-\frac{\sum_{i=1}^{n}(y_i-\widehat y_i^{(m)})^2}
 {\sum_{i=1}^{n}(y_i-\overline y)^2},\qquad
 \overline y=\frac{1}{n}\sum_{i=1}^{n}y_i.
\end{equation}
A definição compara a soma de erros quadráticos com a variação em torno da média \cite{james2023}. Fora do treino, $R^2$ pode ser negativo. Aqui ele resume os casos reunidos, não a média dos coeficientes dos cinco folds; não é uma proporção de variação causada pelas letras.

A diferença pareada e a redução relativa de MAE entre configurações $a$ e $b$ são
\begin{align}
 \Delta_{a-b}&=\frac{1}{n}\sum_{i=1}^{n}
 (|e_i^{(a)}|-|e_i^{(b)}|)=\operatorname{MAE}_a-\operatorname{MAE}_b,\label{eq:delta}\\
 G_{a:b}&=100\frac{\operatorname{MAE}_b-\operatorname{MAE}_a}{\operatorname{MAE}_b}.\label{eq:ganho}
\end{align}
Valor negativo de $\Delta$ favorece $a$; valor positivo de $G$ representa redução percentual do erro. O pareamento exige as mesmas faixas e previsões externas correspondentes. Uma sensibilidade adicional mantém esses ajustes e aplica os pesos originais apenas à avaliação:
\begin{equation}\label{eq:maeponderado}
 \operatorname{MAE}_{m,w}=\frac{\sum_i w_i|e_i^{(m)}|}{\sum_i w_i}.
\end{equation}
Essa escolha não equivale a treinar com pesos e não corrige automaticamente seleção por disponibilidade de letras.

\subsection{Descritores e embeddings de letras}
Descritores explícitos convertem propriedades observáveis de uma letra em números: extensão, diversidade de segmentos e repetição de linhas, por exemplo. Embeddings são vetores densos produzidos por um modelo pré-treinado. Os primeiros têm definições diretamente interpretáveis; os segundos oferecem uma representação aprendida de relações textuais. Nenhuma das duas abordagens garante adequação uniforme entre línguas ou ganho para a resposta de interesse.

\citeonline{parada2024lyrics} analisam características de letras em inglês ao longo do tempo. Seus descritores e sua base derivada do LFM2b \cite{parada2024data} motivaram a trajetória do projeto, mas não são as 21 medidas próprias aqui utilizadas. A distinção evita apresentar uma implementação diferente como reprodução daquele estudo.

O multilingual E5 foi proposto para embeddings multilingues \cite{wang2024}. Neste trabalho, o encoder permanece congelado e não recebe popularidade, artista ou título como entrada adicional. O tokenizer de subpalavras é distinto da segmentação lexical dos descritores. Métodos como SentencePiece permitem tokenização sem exigir segmentação prévia em palavras \cite{kudo2018}; isso não torna comparáveis, por definição, contagens lexicais de diferentes sistemas de escrita. A receita efetivamente aplicada ao corpus é especificada na Seção~\ref{subsec:representacoes}.

\subsection{Validação cruzada aninhada e dependência}
A validação cruzada estima desempenho fora do ajuste. Se a mesma avaliação escolhe hiperparâmetros e anuncia o resultado, a comparação pode ser otimista. No desenho aninhado, as partições internas selecionam hiperparâmetros e as externas avaliam o procedimento escolhido \cite{cawley2010}. A escolha de partições e a estimação de erro de generalização são discutidas em \citeonline{hastie2009}. Em português, \citeonline{izbicki2020} apresentam o aprendizado de máquina sob a perspectiva de risco preditivo, distinguindo avaliação de generalização de simples qualidade de ajuste.

Faixas do mesmo artista, versões relacionadas e letras idênticas podem compartilhar informação. Por isso, a unidade de separação deve reunir as relações identificadas, em vez de sortear cada linha de forma isolada. Os grupos adotados são componentes operacionais; ausência de sobreposição dessas chaves não prova independência estatística universal.

Há dependência adicional entre resultados de CV, pois os conjuntos de treino se sobrepõem. \citeonline{bengio2004} mostram os problemas de estimar a variância ignorando tais correlações. Consequentemente, os cinco folds não são tratados como cinco réplicas independentes, nem usados em um teste $t$ ingênuo. A análise auxiliar dos contrastes condiciona em ajustes já produzidos e tem alcance diferente da incerteza integral de todo o treinamento, conforme explicitado no Apêndice~\ref{ap:incerteza}.

\clearpage\section{Dados e metodologia}\label{sec:metodologia}
\subsection{Catálogo Spotify: origem e análise descritiva}\label{subsec:spotify}
O ponto de partida é o arquivo público \textit{Spotify\_1Million\_Tracks}, disponibilizado no Kaggle pela conta \texttt{amitanshjoshi} \cite{spotifydataset}. O catálogo acessível contém 1.159.764 linhas, 20 colunas e 1.159.764 identificadores de faixa distintos, sem identificadores ausentes ou repetidos. A unidade inicial é a faixa identificada no Spotify; versões e gravações relacionadas permanecem distintas. A fonte não informa um desenho probabilístico de coleta a partir de toda a plataforma. Assim, o catálogo é a população finita acessível deste estudo, não um censo do Spotify.

A resposta corresponde ao campo \texttt{popularity}, mantido separadamente de qualquer contagem de execuções ou visualizações de letras. Conforme a documentação do Spotify, o escore varia de zero a cem e considera execuções e recência \cite{spotifydocs}. A data de medição da resposta no arquivo e a versão exata do conjunto Kaggle não foram confirmadas. O campo \texttt{year} varia de 2000 a 2023 e é utilizado como metadado fornecido; não se presume que identifique o primeiro lançamento da composição.

A popularidade tem média 18,38, desvio-padrão 15,89 e mediana 15. Os quartis inferior e superior são 5 e 29, e o percentil 95 é 49. Há 158.391 faixas com escore zero (13,66\% do catálogo), 499 com escore de pelo menos 80 (0,043\%) e 28 com pelo menos 90. A Figura~\ref{fig:panorama} mostra concentração nos escores baixos e uma cauda superior rara. Trata-se de uma distribuição discreta e assimétrica, sem ajuste ou suposição de normalidade para o alvo.

\figura{spotify_panorama}{Distribuições de popularidade e do campo year no catálogo de 1.159.764 faixas.}{fig:panorama}{Elaboração própria a partir da auditoria do catálogo. Cada barra de popularidade reúne dez valores inteiros, exceto a última, que inclui cem. O campo year não identifica a data de medição da popularidade.}

Agrupando o campo \texttt{year}, 443.801 faixas (38,27\%) estão em 2000--2009, 514.990 (44,40\%) em 2010--2019 e 200.973 (17,33\%) em 2020--2023. As respectivas popularidades médias são 12,66, 19,99 e 26,89. A associação entre ano informado e escore pode refletir recência, composição da fonte e outras diferenças; não estabelece uma tendência longitudinal da mesma música.

O catálogo registra 82 categorias de gênero, com um rótulo por linha e sem ausência nesse campo. Os rótulos são preservados como fornecidos. A Figura~\ref{fig:generos} apresenta as 12 categorias mais frequentes, que juntas correspondem a 21,26\% do catálogo. Black-metal é a maior, com 21.852 faixas (1,88\%), enquanto songwriter é a menor, com 589 (0,05\%). Os rótulos incluem estilos, categorias culturais e descritores de ambiente; não devem ser interpretados como uma taxonomia homogênea ou como participações de mercado na plataforma.

\figura{spotify_generos}{As 12 categorias de gênero mais frequentes no catálogo Spotify utilizado.}{fig:generos}{Elaboração própria a partir das contagens do arquivo original. Os números nas barras são contagens; as outras 70 categorias permanecem na população e na amostra.}

A distribuição por crédito de artista também é desigual. Existem 64.157 strings distintas de crédito, após aparar espaços, com mediana de seis faixas por string e máximo de 4.058. As mil strings mais frequentes reúnem 20,09\% das linhas. Essas strings não são pessoas desambiguadas: podem representar bandas e participações. A dependência potencial entre faixas relacionadas será tratada por agrupamento na validação.

A auditoria identificou 15 créditos de artista e um título ausentes, preservados no catálogo. A busca de termos nos títulos sinalizou 50.267 ocorrências de live/ao vivo, 37.932 de remix e 15.984 de remaster; são marcações heurísticas sobrepostas, não classes exaustivas de versões. Não houve deduplicação por título ou escolha da maior popularidade, pois isso alteraria a unidade e a população analisadas.

As medidas acústicas e musicais já acompanham a fonte; não foram recalculadas a partir de áudio. Danceability, energy, speechiness, acousticness, instrumentalness, liveness e valence são escores entre zero e um; resumem, respectivamente, adequação à dança, intensidade, presença de fala, caráter acústico, ausência de voz, indícios de execução ao vivo e positividade musical. Loudness é expressa em decibéis, tempo em batidas por minuto e duração originalmente em milissegundos \cite{spotifyaudio}. A Tabela~\ref{tab:acusticas} resume essas variáveis no catálogo.

\begin{table}[htbp]\centering\small
\caption{Atributos acústicos e duração no catálogo Spotify.}\label{tab:acusticas}
\begin{tabular}{lrrrrr}\toprule
Variável & Média & Desvio & P25 & Mediana & P75 \\\midrule
danceability & 0,5374 & 0,1845 & 0,4130 & 0,5500 & 0,6770 \\
energy & 0,6397 & 0,2705 & 0,4540 & 0,6940 & 0,8730 \\
loudness & -8,9814 & 5,6822 & -10,8290 & -7,4500 & -5,2760 \\
speechiness & 0,0928 & 0,1268 & 0,0371 & 0,0507 & 0,0890 \\
acousticness & 0,3215 & 0,3550 & 0,0064 & 0,1470 & 0,6400 \\
instrumentalness & 0,2523 & 0,3651 & 0,0000 & 0,0018 & 0,6140 \\
liveness & 0,2230 & 0,2011 & 0,0979 & 0,1340 & 0,2920 \\
valence & 0,4556 & 0,2685 & 0,2260 & 0,4380 & 0,6740 \\
tempo & 121,3771 & 29,7797 & 98,7970 & 121,9310 & 139,9030 \\
Duração (s) & 249,5618 & 149,4262 & 181,0910 & 225,7440 & 286,9135 \\
\bottomrule\end{tabular}
\fonte{Elaboração própria a partir da auditoria. $n=1.159.764$ por variável; valores finitos da fonte, incluindo tempo zero. Desvio com divisor $N-1$.}
\end{table}


Instrumentalness tem mediana 0,0018, mas terceiro quartil 0,6140, indicando heterogeneidade importante. A variável é uma estimativa acústica da fonte, não a classificação textual instrumental do LRCLIB. A duração mediana é 225,74 segundos e o terceiro quartil é 286,91. A auditoria não encontrou ausências nos campos acústicos listados; identificou, porém, 1.199 registros de tempo iguais a zero, incompatíveis com a regra diagnóstica local de valor positivo. Eles permanecem na descrição bruta, sem exclusão retroativa para melhorar resultados.

\FloatBarrier
\subsection{LRCLIB: fonte textual e snapshot local}\label{subsec:lrclib}
O LRCLIB é uma base aberta de letras que recebe contribuições e disponibiliza metadados e textos simples ou sincronizados em formato LRC. Sua documentação descreve identificador de registro, artista, título, álbum, duração e indicação instrumental, além das representações textuais \cite{lrclib}. Um registro da base não equivale necessariamente a uma composição única ou a uma letra pronta para análise: podem existir versões, representações diferentes e ausência de texto utilizável.

Neste estudo foi utilizado o dump público identificado como 12 de outubro de 2026, obtido do domínio de distribuição do LRCLIB e descompactado localmente em SQLite. O resumo da preparação e da busca registra 47,84 GB comprimidos, 261,76 GB descompactados e 33.714.178 registros de metadados percorridos na tabela de faixas. Essas contagens caracterizam o snapshot utilizado; não são o número de letras distintas, de músicas do Spotify correspondentes ou de registros atuais do serviço.

\begin{table}[htbp]\centering\small
\caption{Caracterização documental do snapshot LRCLIB utilizado.}\label{tab:lrclib}
\begin{tabular}{P{.40\linewidth}P{.52\linewidth}}\toprule
Aspecto & Registro disponível \\\midrule
Snapshot & 12 de outubro de 2026 \\
Formato local & Banco SQLite obtido de dump público comprimido \\
Volume comprimido / local & 47,84 GB / 261,76 GB (unidades decimais) \\
Metadados percorridos & 33.714.178 registros de faixas \\
Identificação & Artista, título, álbum e duração; ID próprio do LRCLIB \\
Conteúdo textual & Letra simples, sincronizada, indicação instrumental ou texto ausente \\
Uso no estudo & Recuperar letras das 200.000 faixas sorteadas no Spotify \\
\bottomrule\end{tabular}
\fonte{Documentação LRCLIB e resumo concluído da busca local. O total de metadados provém do checkpoint de indexação, sem nova varredura do banco.}
\end{table}

O uso local tornou possível repetir consultas sobre o mesmo snapshot e examinar candidatos sem depender de respostas mutáveis da API. Índices auxiliares selecionaram metadados dos artistas necessários à busca. Suas contagens não são somadas como se fossem letras novas: os índices são estruturas de acesso ao mesmo banco e podem conter recortes sobrepostos.

Não foi produzido um censo global de idiomas, comprimentos ou disponibilidade de cada tipo de letra no dump. A descrição quantitativa textual apresentada adiante refere-se aos resultados de consulta da amostra Spotify, cujo denominador é conhecido. Essa distinção evita extrapolar a composição dos artistas pesquisados para o LRCLIB inteiro. A existência de conteúdo público tampouco garante correspondência correta ou cobertura uniforme por língua, gênero e popularidade.

O LFM2b teve um papel anterior de motivação e integração de descritores, conforme a base de \citeonline{parada2024data}. Seus dados e resultados não foram transportados para preencher a coorte atual. Contagens Last.fm e visualizações Genius permanecem respostas diferentes do escore Spotify. A base principal desta análise é construída pela sequência Spotify--amostragem--LRCLIB descrita a seguir.

\FloatBarrier
\subsection{Seleção probabilística de 200 mil faixas}\label{subsec:amostragem}
A busca de letras e a construção das representações exigem recursos de armazenamento e processamento. Fixou-se, por isso, um orçamento de 200.000 faixas antes da recuperação textual. Esse total é uma decisão operacional; não decorre de cálculo de poder, tamanho ótimo ou garantia genérica de erro. A estratificação utiliza informações disponíveis no catálogo e permite preservar probabilidades e pesos conhecidos \cite{cochran1977,statcan2010}.

As 499 faixas com popularidade de pelo menos 80 foram incluídas integralmente, dentro das 200.000, para preservar a cauda rara do catálogo. As demais foram estratificadas pelo cruzamento de gênero, três períodos do campo year (2000--2009, 2010--2019 e 2020--2023) e cinco bandas de popularidade (0--19, 20--39, 40--59, 60--69 e 70--79). O desenho resultou em 981 estratos não vazios, dos quais 62 foram integrais, incluindo o grupo de alta popularidade.

Para explicitar a alocação, seja $C_0=499$ o tamanho do grupo incluído integralmente e $\mathcal H_0$ o conjunto dos outros estratos. Define-se uma reserva $b_h=\min(2,N_h)$, uma capacidade residual $c_h=N_h-b_h$ e o orçamento restante
\begin{equation}\label{eq:orcamento}
 R=200.000-C_0-\sum_{h\in\mathcal H_0}b_h,\qquad
 q_h=R\frac{c_h}{\sum_{\ell\in\mathcal H_0}c_\ell}.
\end{equation}
A alocação inicial é $n_h=b_h+\lfloor q_h\rfloor$. As vagas ainda disponíveis são distribuídas pelos maiores restos $q_h-\lfloor q_h\rfloor$, com desempate determinístico pelo identificador de estrato. Depois, realiza-se sorteio aleatório simples sem reposição em cada estrato não integral. A distribuição é aproximadamente proporcional à capacidade residual, não uma alocação de Neyman ou otimização demonstrada para o MAE.

O sorteio utilizou gerador PCG64 e semente 20261031, com ordenação fixa de estratos e identificadores. As probabilidades e pesos seguem a Equação~\ref{eq:inclusao}. A amostra contém exatamente 200.000 IDs únicos, inclui as 499 faixas da cauda e conserva a soma dos pesos igual a 1.159.764. Nenhuma faixa foi substituída por possuir letra mais fácil de localizar. Gêneros e anos originais, pesos e estratos permanecem disponíveis para diagnóstico de cobertura.

\subsection{Integração por metadados e aceitação de letras}\label{subsec:integracao}
Somente depois do sorteio foram consultadas as correspondências no LRCLIB. A recuperação utiliza artista, título e duração da faixa Spotify, com referências separadas ao registro e ao conteúdo textual encontrado. Popularidade, gênero, peso amostral e instrumentalness não participam da decisão de aceitar uma letra, evitando que a resposta ou um preditor acústico sejam usados como evidência de identidade textual.

A normalização cria chaves adicionais e preserva os campos originais. Nas consultas finais, uniformizam-se Unicode, caixa, espaços e sinais tipográficos usuais, preservando acentos e nomes formados apenas por símbolos. Participações são reconhecidas quando explicitadas por feat., ft. ou featuring. Não se decompõem automaticamente nomes de bandas e créditos por qualquer vírgula ou sinal \&. Qualificadores contextuais, como live, remix, remaster e versões acústicas, são conservados para impedir que a simplificação do título iguale gravações incompatíveis.

Os candidatos precisam ter artista e título compatíveis, incluindo as informações de colaboração e versão disponíveis. Quando ambas as durações são válidas, exige-se diferença de no máximo dez segundos; ausência de duração não comprova identidade. Não se aceita uma correspondência apenas por título isolado, substring ou maior similaridade. Se candidatos elegíveis contêm textos substantivamente diferentes ou conflitos de versão, a decisão permanece ambígua.

Letras sincronizadas são convertidas somente nos formatos LRC suportados, removendo marcadores temporais e metadados sem apagar conteúdo, ordem e repetições. Quando texto simples e sincronizado estão disponíveis, divergências impedem novas aceitações. Diferenças puramente tipográficas têm equivalência restrita; o procedimento não escolhe a letra mais conveniente. Uma primeira recuperação já havia aceitado 77.391 faixas; suas aceitações e classificações instrumentais foram preservadas na recuperação final, que acrescentou 144 aceitas. Isso significa que parte das decisões foi herdada, sem uma revalidação humana universal.

\begin{table}[htbp]
\centering\small
\caption{Destino das faixas da amostra após a recuperação textual.}
\label{tab:cobertura}
\begin{tabular}{lrr}
\toprule
Situação & Faixas & Percentual da amostra \\
\midrule
Aceitas & 77.535 & 38,7675\% \\
Não encontradas & 86.610 & 43,3050\% \\
Ambíguas & 27.085 & 13,5425\% \\
Instrumentais & 8.759 & 4,3795\% \\
Sem texto & 9 & 0,0045\% \\
Metadados insuficientes & 2 & 0,0010\% \\
\bottomrule
\end{tabular}
\fonte{Elaboração própria a partir dos agregados auditados do estudo.}
\nota{Total de 200.000 registros; zero erro técnico ou pendência. As não aceitas permanecem na base de cobertura.}
\end{table}


O resultado final é de 77.535 letras aceitas, 38,77\% da amostra. As outras 122.465 faixas permanecem registradas nos diagnósticos de cobertura, e não foram substituídas. Não encontrada significa ausência de candidato aceito sob as chaves, critérios e snapshot utilizados, não inexistência universal de letra. Ambiguidade representa evidência insuficiente para escolher entre candidatos; instrumental é uma classificação da fonte textual, distinta do preditor instrumentalness. Não restaram erros técnicos ou consultas pendentes.

Integridade de vínculo e texto foi auditada, mas não se estimou uma taxa global de precisão do matching por revisão humana probabilística. Identificadores e hashes verificam consistência técnica, não a pertinência semântica de toda letra à gravação. Homônimos, pseudônimos e decisões herdadas são limitações da integração, mantidas na interpretação dos resultados.

\FloatBarrier
\subsection{Da amostra sorteada ao domínio observado}\label{subsec:selecao}
A Tabela~\ref{tab:descritivas} compara catálogo, amostra e aceitas. A média da amostra ponderada, 18,3899, aproxima-se da média do catálogo, 18,3831, enquanto a média ponderada entre aceitas é 23,7482. Sem pesos, a média das aceitas é 23,9821 e sua mediana é 22, contra 15 no catálogo. A Figura~\ref{fig:distribuicao} mostra o deslocamento da distribuição observada após a etapa textual.

\begin{table}[htbp]
\centering\small
\caption{Popularidade observada: catálogo, amostra e domínio com letras aceitas.}
\label{tab:descritivas}
\begin{tabular}{lrrrrr}
\toprule
Universo / avaliação & n & Média & P25 & Mediana & P75 \\
\midrule
Catálogo & 1.159.764 & 18,3831 & 5 & 15 & 29 \\
Amostra, sem pesos & 200.000 & 18,6752 & 5 & 15 & 29 \\
Amostra, ponderada & 200.000 & 18,3899 & 5 & 15 & 29 \\
Aceitas, sem pesos & 77.535 & 23,9821 & 12 & 22 & 34 \\
Aceitas, ponderadas & 77.535 & 23,7482 & 12 & 22 & 34 \\
\bottomrule
\end{tabular}
\fonte{Elaboração própria a partir dos agregados auditados do estudo.}
\nota{Quantis pela inversa da distribuição empírica, ponderada quando indicado. n conta faixas, não soma de pesos.}
\end{table}

\figura{distribuicao_popularidade}{Distribuição acumulada da popularidade antes e depois da recuperação de letras.}{fig:distribuicao}{Elaboração própria a partir das frequências agregadas. Amostra ponderada pelos pesos de desenho; aceitas com pesos iguais por faixa.}

A cobertura global ponderada é 38,7436\%, próxima da proporção sem pesos. Entretanto, a Figura~\ref{fig:cobertura} evidencia que essa proximidade não elimina seleção por popularidade. Das 499 faixas incluídas integralmente na banda de pelo menos 80, apenas 128 tiveram letras aceitas (25,65\%). Há ainda 102 estratos sem aceitas e 348 com menos de cinco aceitas. Um peso original positivo não cria informação textual nos estratos não observados.

\figura{cobertura_bandas}{Proporção de letras aceitas por banda de popularidade na amostra.}{fig:cobertura}{Elaboração própria a partir da recuperação local. Cada denominador é o número de faixas sorteadas na banda; proporções sem pesos.}

Um diagnóstico exploratório anterior comparou 104 critérios marginais com tolerâncias explícitas: um ponto para a média de popularidade; 0,10 desvio-padrão do catálogo para outras médias; e o menor entre um ponto percentual e 20\% da proporção do catálogo para categorias. A análise por linearização sob o desenho utilizou equivalência de Wald (TOST a 5\%) e controle simultâneo de diferenças por Bonferroni. Foram 21 critérios equivalentes, 74 diferenças individuais e nove inconclusivos; 72 diferenças tiveram apoio no controle simultâneo. O resultado não sustenta equivalência global das aceitas ao catálogo. Esses diagnósticos são posteriores a explorações do projeto e não constituem confirmação independente.

A conclusão sobre inferência é delimitada: a seleção probabilística inicial permite descrição e estimação sob o desenho para o catálogo definido; a posterior recuperação textual não preserva automaticamente esse alcance. Não se demonstrou que as 77.535 aceitas sejam uma amostra representativa de todo o catálogo ou de todas as músicas com letras. Por isso, o estimando principal da modelagem é o desempenho preditivo médio nas avaliações externas da distribuição empírica desse domínio observado, com pesos iguais por faixa. Os pesos originais são preservados para descrição e sensibilidade, sem serem interpretados como correção automática da disponibilidade textual.

O domínio permaneceu multilingue: 56.908 faixas foram classificadas no status inglês, 20.447 no status não inglês e 180 como incertas. O status combina código e confiança de um detector automático, não uma verificação humana. Houve 77.051 entradas textuais normalizadas distintas, de modo que 484 vínculos excedentes compartilham conteúdo. A análise conserva as faixas e considera texto igual na formação de grupos de validação.

\FloatBarrier
\subsection{Preditores e representações utilizadas}\label{subsec:representacoes}
O Baseline contém 15 campos: 11 numéricos (year, danceability, energy, loudness, speechiness, acousticness, instrumentalness, liveness, valence, tempo e duração) e quatro categóricos (genre, key, mode e time\_signature). Key representa a classe de altura musical; mode codifica modo maior/menor; time\_signature informa a assinatura de tempo \cite{spotifyaudio}. Gênero é a categoria da fonte, e não uma classificação construída a partir da letra. Nomes, IDs, idioma, chaves de grupos, pesos e informações de matching não entram como preditores.

\begin{table}[htbp]\centering\small
\caption{Configurações preditivas comparadas nas mesmas faixas e folds.}\label{tab:blocos}
\begin{tabular}{lP{.65\linewidth}}\toprule
Configuração & Conjunto de características antes da codificação categórica \\\midrule
Baseline & 15 campos acústicos, musicais e de metadados \\
Tradicional & Baseline + 21 descritores lexicais e estruturais próprios \\
Profundo & Baseline + 384 coordenadas do embedding multilingue \\
Combinado & Baseline + 21 descritores + 384 coordenadas \\
\bottomrule\end{tabular}
\fonte{Elaboração própria. A quantidade final de colunas depende das categorias observadas em cada treino. Todos os regressores são Ridge.}
\end{table}

Descritores e embeddings recebem o mesmo texto aceito, com normalização Unicode NFC, uniformização de quebras de linha e remoção de espaços finais e linhas vazias externas. Não se removem palavras, repetições ou pontuação. Os 21 descritores resumem extensão, diversidade de segmentos, entropia, estrutura de linhas e composição gráfica. São definidos integralmente no Apêndice~\ref{ap:descritores}; não constituem reprodução dos descritores LFM2b nem medidas de sentimento validadas para todas as línguas. Segmentos lexicais são sequências de letras Unicode com apóstrofo ou hífen interno, e não palavras linguísticas universais.

O encoder é o \textit{multilingual-e5-small}, revisão congelada \texttt{614241f6}, com vetor de 384 dimensões, execução em CPU e ponto flutuante de 32 bits, sem ajuste com popularidade \cite{wang2024,e5card}. O texto é tokenizado integralmente e dividido em janelas com sobreposição de 32 tokens de conteúdo. Reservam-se prefixo e tokens especiais no limite de 512 posições; não se descarta silenciosamente o final de letras longas.

Para uma janela $j$, sejam $\bm h_{jt}$ os estados dos tokens e $a_{jt}\in\{0,1\}$ sua máscara de atenção. O pooling e a agregação são
\begin{equation}\label{eq:embedding}
 \bm v_j=\frac{\sum_t a_{jt}\bm h_{jt}}{\sum_t a_{jt}},\quad
 \bm u_j=\frac{\bm v_j}{\|\bm v_j\|_2},\quad
 \bm z=\frac{\sum_j c_j\bm u_j}{\|\sum_j c_j\bm u_j\|_2},
\end{equation}
em que $c_j$ é o número de tokens de conteúdo novos na janela. Padding é excluído; prefixo e tokens especiais válidos participam da média. O pooling mascarado e a normalização L2 seguem a orientação do model card; a agregação de janelas é uma adaptação deste estudo, não garantia de preservação integral de semântica.

Há um desvio de entrada que precisa qualificar os resultados. O model card recomenda o prefixo \texttt{query: } para uso como características. A execução tokenizou prefixo e conteúdo separadamente, introduzindo um token adicional de espaço em comparação com a concatenação textual da receita oficial. Em quatro frases artificiais, as duas receitas produziram vetores próximos, mas diferentes, com cossenos entre 0,99664 e 0,99943. Não foi medido o efeito dessa diferença no MAE do corpus. A representação avaliada é, portanto, uma variante operacional documentada do E5; não se alega reprodução canônica exata.

Foram produzidos 77.051 vetores, associados às 77.535 faixas por referências de conteúdo. A auditoria verificou dimensão, finitude e norma dos vetores; uma conferência limitada posterior reproduziu os 21 descritores em oito textos, com diferença máxima zero. Esses controles verificam implementação e alinhamento, sem certificar adequação semântica ou uniformidade entre línguas. As representações e a população foram preservadas após conhecer os resultados.

\subsection{Validação, ajuste e política de comparação}\label{subsec:validacao}
A separação dos dados utiliza componentes conexos de relações identificadas por artista, composição operacional, versão e texto igual. Se uma relação liga duas faixas, elas e suas ligações transitivas permanecem no mesmo componente. Foram formados 21.893 componentes, com máximo de 186 faixas. As chaves são operacionais e podem deixar de reconhecer aliases ou versões; não são identificadores universais de composição.

Aplicou-se validação cruzada agrupada aninhada 5$\times$3, com os mesmos componentes e partições nas quatro configurações. Cada teste externo contém 15.507 faixas e seu treino 62.028. Dentro do treino externo, cada avaliação interna contém 20.676 faixas e seu treino 41.352. A distribuição equilibra quantidades de faixas sem usar popularidade para formar ou redesenhar os grupos.

Em cada treino, as características numéricas passam por imputação de mediana e padronização conforme a Equação~\ref{eq:padronizacao}. Nos 15 campos Baseline da coorte final não houve ausência ou invalidade registrada. As categorias são codificadas por indicadores binários, com vocabulário aprendido no treino e zeros para categorias novas na avaliação. Nenhuma transformação aprende parâmetros no teste. Não há seleção global de variáveis ou redução por componentes principais.

Para cada configuração e treino externo, a penalização é selecionada na grade
$\mathcal A=\{0{,}1;1;10;100;1.000;10.000\}$. Se $V_{k\ell}$ indica a avaliação interna $\ell$ do treino externo $k$, escolhe-se
\begin{equation}\label{eq:selecaoalpha}
 \widehat\alpha_{k,m}=\underset{\alpha\in\mathcal A}{\operatorname{argmin}}
 \frac{\sum_{\ell=1}^{3}\sum_{i\in V_{k\ell}}
 |y_i-\widehat f^{(m)}_{k\ell,\alpha}(\bm x_i)|}
 {\sum_{\ell=1}^{3}|V_{k\ell}|}.
\end{equation}
Empates exatos favorecem o maior $\alpha$. O pré-processamento e o Ridge são então reajustados sobre todo o treino externo, com a penalização escolhida, para prever seu teste. Reunidas as cinco avaliações externas, cada faixa tem uma previsão por configuração. MAE, RMSE e $R^2$ seguem as Equações~\ref{eq:mae}--\ref{eq:r2}; os ajustes e as métricas principais dão peso igual a cada faixa, sem equilibrar artistas por meio de pesos artificiais.

Duas referências adicionais predizem a média e a mediana do respectivo treino externo. A sensibilidade da Equação~\ref{eq:maeponderado} utiliza os pesos originais somente nas mesmas previsões externas, sem retreinar ou selecionar novamente os modelos. Os ajustes foram realizados em Python e scikit-learn \cite{pedregosa2011}, com solver LSQR, tolerância $10^{-5}$ e limite de 10.000 iterações. Os 360 ajustes internos Ridge não registraram avisos ou alcance desse limite.

A comparação atual sucede explorações anteriores do projeto. Um teste histórico foi consultado em 14 de outubro de 2026, havendo vínculos de artista e composição com o domínio recente. O agrupamento atual evita as sobreposições operacionais especificadas, mas um novo sorteio não apaga essa exposição. A avaliação é apresentada como estudo empírico com essa história, sem alegar confirmação independente. Os cinco folds tampouco fornecem cinco réplicas independentes para inferência.

\clearpage\section{Resultados e discussão}\label{sec:resultados}
\subsection{Desempenho global e ganho incremental}
A Tabela~\ref{tab:metricas} reúne as previsões externas das mesmas 77.535 faixas. Todas as configurações Ridge superaram as referências de média e mediana do treino. Entre as configurações, o Combinado apresentou o menor MAE (8,1151), seguido de Profundo (8,1285), Tradicional (8,2957) e Baseline (8,3525). RMSE e $R^2$ mantiveram a mesma ordenação. O resultado mostra redução de erro na avaliação realizada, sem atribuir às letras causalidade sobre o escore.

\begin{table}[htbp]
\centering\small
\caption{Desempenho nas previsões externas reunidas ($n=77.535$).}
\label{tab:metricas}
\begin{tabular}{lrrrr}
\toprule
Configuração & MAE & RMSE & $R^2$ & MAE ponderado \\
\midrule
Baseline & 8,3525 & 10,8255 & 0,5087 & 8,2194 \\
Tradicional & 8,2957 & 10,7538 & 0,5152 & 8,1652 \\
Profundo & 8,1285 & 10,5540 & 0,5330 & 8,0032 \\
Combinado & 8,1151 & 10,5381 & 0,5344 & 7,9908 \\
Média do treino & 12,5621 & 15,4463 & -0,0003 & 12,4181 \\
Mediana do treino & 12,4786 & 15,5708 & -0,0165 & 12,3213 \\
\bottomrule
\end{tabular}
\fonte{Elaboração própria a partir dos agregados auditados do estudo.}
\nota{Pesos iguais no ajuste e na avaliação principal. Última coluna: sensibilidade com pesos originais nas mesmas previsões; não houve retreino ponderado.}
\end{table}


O acréscimo dos 21 descritores ao Baseline reduziu o MAE em 0,0568 ponto (0,68\%). O acréscimo dos embeddings reduziu-o em 0,2240 ponto (2,68\%). Reunindo ambas as representações, a redução frente ao Baseline foi 0,2374 ponto (2,84\%). Assim, a maior parcela do ganho conjunto já aparece na configuração com embeddings. Adicionar os descritores ao Profundo produziu redução de apenas 0,0134 ponto (0,16\%). A dimensão do ganho importa além de sua direção: a melhora média é modesta diante de um erro absoluto ainda próximo de oito pontos.

A ponderação pelos pesos originais de amostragem, aplicada somente às mesmas previsões, preservou a ordenação. O MAE ponderado passou de 8,2194 no Baseline a 7,9908 no Combinado, redução de 2,78\%. Essa concordância indica que o sentido da comparação não depende apenas da contribuição relativa dos estratos na avaliação sem pesos. Não demonstra que os pesos corrijam a seleção de letras ou que o desempenho ponderado se generalize ao catálogo.

O Combinado superou o Baseline nos cinco testes externos; as diferenças de MAE variaram de $-0,2724$ a $-0,2070$. Retirar um componente de cada vez apenas da avaliação, mantendo os modelos fixos, não inverteu o sinal dos contrastes examinados. Esses resultados descrevem estabilidade nos casos e divisões existentes, mas não são cinco réplicas independentes, validação com reajuste a cada retirada ou intervalos de confiança.

Em seis dos vinte ajustes externos Ridge, a penalização selecionada foi o menor valor da grade, $\alpha=0,1$. Isso não demonstra erro de otimização ou que uma grade ampliada necessariamente melhoraria desempenho. A comparação se refere ao procedimento especificado; modificar a grade depois de observar seus testes externos criaria outra etapa de desenvolvimento. Os resultados atuais foram preservados.

\subsection{Incerteza dos contrastes e alcance da evidência}\label{subsec:incerteza}
As diferenças calculadas nas previsões existentes são conhecidas exatamente para esses casos. Quantificar sua variação em avaliações hipotéticas exige um alvo e hipóteses adicionais. A análise auxiliar do Apêndice~\ref{ap:incerteza} utiliza os componentes para permitir dependência interna e condiciona separadamente no par de modelos de cada fold. Os limites da Tabela~\ref{tab:incerteza} têm nível nominal de 95\% nesse modelo auxiliar, com ajuste Bonferroni; não representam a incerteza integral de seleção, treinamento e construção de dados.

\begin{table}[htbp]
\centering\small
\caption{Limites auxiliares condicionais para os contrastes de interesse.}
\label{tab:incerteza}
\begin{tabular}{lrrrr}
\toprule
Contraste & Avaliação & Delta & Inferior & Superior \\
\midrule
C - B & Iguais & -0,23736 & -0,31400 & -0,16073 \\
C - P & Iguais & -0,01340 & -0,03027 & 0,00347 \\
C - B & Originais & -0,22861 & -0,30454 & -0,15268 \\
C - P & Originais & -0,01234 & -0,02919 & 0,00450 \\
\bottomrule
\end{tabular}
\fonte{Elaboração própria a partir dos agregados auditados do estudo.}
\nota{C: Combinado; B: Baseline; P: Profundo. 95\% nominais sob modelo auxiliar de componentes, com Bonferroni por família de dez marginais. Não são ICs da pipeline integral.}
\end{table}


Para Combinado--Baseline, os limites mantêm valores negativos, entre $-0,3140$ e $-0,1607$ ponto. Para Combinado--Profundo, abrangem zero, de $-0,0303$ a $0,0035$. A sensibilidade ponderada leva à mesma interpretação. O ganho adicional dos descritores além dos embeddings permanece inconclusivo nessa análise; isso não demonstra equivalência entre configurações. Também não se declara significância integral do ganho conjunto, pois o alvo condicional não abrange toda a variação do procedimento.

Essa distinção é compatível com a discussão de \citeonline{bates2024}: erro de um modelo ajustado, desempenho médio de um algoritmo e variabilidade da CV não são a mesma quantidade. Um bootstrap das perdas por componente com previsões congeladas variaria a composição da avaliação, mas não repetiria a escolha de hiperparâmetros nem o treinamento. A grande quantidade de componentes, por si só, não resolve a sobreposição entre treinos ou transforma grupos operacionais em uma amostra independente de composições.

\subsection{Popularidade elevada e subestimação}
A Figura~\ref{fig:bandaserro} mostra que a média global encobre erros diferentes ao longo do alvo. Nas 128 faixas aceitas com popularidade de pelo menos 80, o Combinado teve MAE de 29,3711 pontos, muito acima do global. Essas faixas pertencem a 96 componentes e têm popularidade observada entre 80 e 96: 116 estão em 80--84, oito em 85--89 e quatro em 90--100. Portanto, a dificuldade não se restringe a poucos valores próximos de cem.

\figura{erro_popularidade}{MAE do Baseline e do Combinado por banda de popularidade observada.}{fig:bandaserro}{Elaboração própria a partir das previsões externas existentes. As bandas são definidas pelo alvo observado; contagens indicadas na figura.}

Todas as 128 faixas foram subestimadas pelas quatro configurações. No Combinado, a popularidade média observada foi 81,8125, enquanto a previsão média foi 52,4414; a maior previsão nessa cauda foi 77,9977. Definindo o desvio assinado de diagnóstico por $r_i=\widehat y_i-y_i=-e_i$, a mediana foi $-28,7755$, com primeiro e terceiro quartis de $-37,2444$ e $-22,4381$. Houve 101 subestimações de pelo menos vinte pontos e 22 de pelo menos quarenta pontos, evidenciando um padrão amplo na cauda aceita.

Apesar desses erros elevados, a cauda contém apenas 0,1651\% das faixas avaliadas e produz 0,5975\% do erro absoluto total do Combinado. Sua contribuição aditiva ao MAE global é
\begin{equation}\label{eq:contribuicaocauda}
 \frac{1}{77.535}\sum_{i:y_i\ge80}|e_i|=0,04849.
\end{equation}
Para o erro quadrático total, a participação é 1,4591\%. O baixo peso numérico da cauda permite que um bom resultado médio relativo conviva com desempenho fraco nos escores mais altos. Isso limita aplicações voltadas a identificar músicas muito populares.

A compressão das previsões é compatível com regularização e informação insuficiente para distinguir a cauda, mas o diagnóstico não identifica causalmente seu mecanismo. Condicionar nos maiores valores observados pode gerar erros negativos mesmo sem falha de código. Essa análise também difere de estudar calibração condicionada na previsão, $E(Y\mid\widehat Y)$. No domínio inteiro, o Combinado não produziu nenhuma previsão de pelo menos 80.

\FloatBarrier
\subsection{Previsões negativas e restrição ao suporte}
O Combinado produziu 615 previsões negativas, em 317 componentes, com mínimo de $-32,7066$ e mediana de $-1,7873$. Nessas faixas, a popularidade observada tem média 5,3122 e mediana três; 133 respostas são zero e apenas uma é de pelo menos quarenta. A extrapolação se concentra na borda inferior. A regressão Ridge da Equação~\ref{eq:ridge} não impõe restrição de suporte, de modo que valores negativos não constituem, isoladamente, evidência de dados corrompidos.

Como sensibilidade separada, aplicou-se a projeção
\begin{equation}\label{eq:clipping}
 \widetilde y_i=\min\{100,\max\{0,\widehat y_i\}\}.
\end{equation}
Para $y_i\in[0,100]$, ela não aumenta o erro absoluto nem o quadrático de nenhuma observação, pois aproxima uma previsão externa ao intervalo admissível. A melhora decorrente não representa nova informação das letras. A Tabela~\ref{tab:clipping} apresenta essa análise sem substituir o resultado principal.

\begin{table}[htbp]
\centering\small
\caption{Projeção de suporte como sensibilidade pós-resultado.}
\label{tab:clipping}
\begin{tabular}{lrrrr}
\toprule
Configuração & Abaixo de 0 & Acima de 100 & MAE original & MAE projetado \\
\midrule
Baseline & 537 & 0 & 8,35248 & 8,33781 \\
Tradicional & 584 & 1 & 8,29569 & 8,28013 \\
Profundo & 610 & 0 & 8,12852 & 8,11055 \\
Combinado & 615 & 0 & 8,11512 & 8,09701 \\
\bottomrule
\end{tabular}
\fonte{Elaboração própria a partir dos agregados auditados do estudo.}
\nota{Projeção fixa em [0,100], sem substituir as previsões ou métricas principais.}
\end{table}


A ordenação permanece a mesma e o MAE do Combinado passa a 8,0970, redução de 0,0181 ponto. Nenhuma das 128 previsões da cauda é alterada, porque elas já estavam em $[0,100]$. O clipping reduz extrapolações de suporte, mas não resolve a subestimação das faixas de alta popularidade. Como análise posterior, não foi incorporado retroativamente à regra principal de avaliação.

\subsection{Idiomas, heterogeneidade e limitações}
A Tabela~\ref{tab:idiomas} evidencia heterogeneidade linguística. O ganho Combinado--Baseline é favorável nos status inglês e não inglês, mas desfavorável nos 180 casos incertos. Esse status pequeno não permite uma conclusão geral sobre uma língua específica. Os grupos são estimados automaticamente, e nenhuma faixa foi excluída para melhorar a comparação.

\begin{table}[htbp]
\centering\small
\caption{Desempenho por status do idioma estimado.}
\label{tab:idiomas}
\begin{tabular}{lrrrr}
\toprule
Status & n & Baseline & Profundo & Combinado \\
\midrule
Inglês estimado & 56.908 & 7,6840 & 7,5418 & 7,5278 \\
Não inglês estimado & 20.447 & 10,2036 & 9,7491 & 9,7343 \\
Idioma incerto & 180 & 9,4360 & 9,5398 & 9,8564 \\
\bottomrule
\end{tabular}
\fonte{Elaboração própria a partir dos agregados auditados do estudo.}
\nota{Idioma automático não é rótulo humano. Diagnóstico exploratório sem excluir casos.}
\end{table}


Entre os 76 gêneros com métricas publicadas pelos limiares descritivos de pelo menos cem faixas e vinte componentes, 59 favoreceram o Combinado e 17 não. Essas comparações são exploratórias e podem compartilhar componentes entre categorias; não representam uma bateria confirmatória de testes por gênero. O ganho médio não autoriza prometer benefício para todos os idiomas ou estilos.

O alcance das conclusões combina limitações de fonte, integração e avaliação. O catálogo não cobre probabilisticamente toda a plataforma, a data do escore é desconhecida e a disponibilidade de letras é seletiva. As regras de matching preservam ambiguidades, mas não tiveram precisão global estimada por revisão humana probabilística. O agrupamento reduz vazamentos identificados, sem resolver toda identidade musical ou exposição histórica. Por fim, o encoder avaliado é a variante operacional descrita, e não uma garantia abstrata de desempenho de qualquer embedding E5. Esses limites qualificam a evidência obtida sem exigir que resultados de outras populações sejam misturados aos atuais.

\clearpage\section{Considerações finais}\label{sec:conclusao}
Este trabalho examinou se representações de letras acrescentam informação preditiva a atributos acústicos e metadados na estimação da popularidade Spotify observada. A construção dos dados foi parte central da análise: de um catálogo acessível com 1.159.764 faixas, selecionaram-se 200.000 por amostragem estratificada e recuperaram-se letras aceitas para 77.535. A descrição mostrou que esse domínio difere do catálogo em popularidade e cobertura. A seleção probabilística inicial não autoriza tratar automaticamente o subconjunto textual como representativo da população de origem.

Na comparação Ridge com validação agrupada aninhada comum, o MAE passou de 8,3525 no Baseline a 8,1151 no Combinado, redução de 2,84\%. A configuração com embeddings apresentou a maior parcela desse ganho; o acréscimo posterior dos 21 descritores produziu redução menor, de 0,0134 ponto, inconclusiva na análise auxiliar de incerteza. A sensibilidade com pesos originais apenas na avaliação preservou a ordenação, sem corrigir automaticamente a disponibilidade de letras.

O resultado médio não se reproduziu uniformemente nos subgrupos. Todas as 128 faixas aceitas de popularidade elevada foram subestimadas, e os casos de idioma incerto não apresentaram ganho conjunto. Previsões negativas decorreram de um estimador sem restrição de suporte; a sensibilidade por clipping reduziu ligeiramente o MAE, mas não melhorou a cauda. Assim, a evidência sustenta contribuição textual modesta no procedimento e domínio estudados, sem demonstrar capacidade prospectiva de prever sucessos ou efeito causal das letras.

Como continuidade, são pertinentes a confirmação de metadados temporais da fonte, uma revisão humana planejada da qualidade de correspondência e a avaliação independente de uma representação com prefixo corrigido. Métodos para ampliar o alcance inferencial dependeriam de hipóteses explícitas sobre seleção textual e suporte, além de validação apropriada. Essas possibilidades constituem novos estudos delimitados, sem substituir ou reinterpretar retrospectivamente a comparação aqui apresentada.

\clearpage\phantomsection\addcontentsline{toc}{section}{Referências}
\markboth{Referências}{}
\bibliographystyle{abntex2-alf}
\begingroup\small\singlespacing\raggedright
\interlinepenalty=10000
\bibliography{ref}
\endgroup
\clearpage\appendix
\section{Definições dos descritores textuais}\label{ap:descritores}
As medidas são calculadas no texto normalizado comum. $C$ conta caracteres não brancos; $L$, letras Unicode; $T$, segmentos lexicais; $U$, tipos distintos; e $B$, linhas não vazias. Segmentos são sequências de letras com apóstrofo ou hífen interno, após casefold. Denominadores nulos recebem zero, com indicador de ausência de segmentos.

\begingroup\footnotesize\singlespacing
\begin{longtable}{rP{.31\linewidth}P{.56\linewidth}}
\caption{Definições operacionais dos 21 descritores de letras.}\label{tab:descritores}\\
\toprule & Descritor & Definição \\\midrule\endfirsthead
\toprule & Descritor & Definição (continuação) \\\midrule\endhead
\bottomrule\endfoot
1 & Caracteres totais & Pontos de código, incluindo espaços e quebras internas.\\
2 & Caracteres não brancos & Contagem $C$.\\
3 & Linhas não vazias & Contagem $B$.\\
4 & Estrofes operacionais & Blocos não vazios separados por uma ou mais linhas vazias.\\
5 & Segmentos lexicais & Contagem $T$.\\
6 & Tipos lexicais & Contagem $U$, após casefold.\\
7 & Comprimento médio de segmento & Média de caracteres dos segmentos após casefold.\\
8 & Dispersão de comprimento de segmento & Desvio com divisor $T$.\\
9 & Razão tipo/segmento & $U/T$.\\
10 & Razão de tipos únicos em frequência & Número de tipos com uma ocorrência dividido por $U$.\\
11 & Entropia de segmentos & $-\sum_{t=1}^{U}p_t\log_2p_t$, com frequências relativas $p_t$.\\
12 & Comprimento médio de linha & Média de caracteres não brancos nas linhas não vazias.\\
13 & Dispersão de comprimento de linha & Desvio com divisor $B$.\\
14 & Fração de linhas repetidas & $(B-\text{número de linhas distintas})/B$, após casefold e uniformização de espaços.\\
15 & Fração de letras & $L/C$, com categoria Unicode L.\\
16 & Fração de dígitos & Dígitos decimais Unicode divididos por $C$.\\
17 & Fração de pontuação & Caracteres de categoria Unicode P divididos por $C$.\\
18 & Fração de maiúsculas & Letras maiúsculas divididas por $L$.\\
19 & Linhas com interrogação & Fração das $B$ linhas com interrogação simples ou de largura completa.\\
20 & Linhas com exclamação & Fração das $B$ linhas com exclamação simples ou de largura completa.\\
21 & Ausência de segmentos & Indicador de $T=0$.\\
\end{longtable}
\endgroup

A entropia segue a definição discreta em bits de \citeonline{cover2006}, aplicada aqui aos segmentos. Ela e a razão tipo/segmento dependem da extensão e segmentação, sem interpretação linguística universal. O indicador 21 foi constante zero e é centrado para zero na padronização.

\clearpage\section{Limites auxiliares para os contrastes}\label{ap:incerteza}
\setcounter{equation}{0}
\renewcommand{\theequation}{\thesection.\arabic{equation}}
Esta análise é exploratória e condicionada às previsões produzidas. Para um contraste $a-b$, defina $d_i=|e_i^{(a)}|-|e_i^{(b)}|$. Em um componente $g$ do teste externo $k$, sejam $D_g=\sum_{i\in g}d_i$ e $N_g$ o número de faixas. O contraste no fold e seu erro-padrão auxiliar são
\begin{align}
 \widehat\delta_k&=\frac{\sum_{g\in k}D_g}{\sum_{g\in k}N_g},\label{eq:deltafold}\\
 \widehat{\operatorname{SE}}_k^2&=
 \frac{G_k}{G_k-1}\frac{\sum_{g\in k}(D_g-\widehat\delta_k N_g)^2}
 {(\sum_{g\in k}N_g)^2},\label{eq:sefold}
\end{align}
em que $G_k$ é o número de componentes do teste. A razão preserva peso igual por faixa, em vez de peso igual por componente. Na sensibilidade ponderada, usam-se $D_g=\sum_{i\in g}w_i d_i$ e $N_g=\sum_{i\in g}w_i$.

A linearização admite dependência interna ao componente e pressupõe independência, momentos adequados e ausência de grupos dominantes entre componentes, sob um modelo auxiliar de avaliações hipotéticas \cite{cameron2015}. O sorteio original de faixas não comprova essas hipóteses para os grupos formados posteriormente; relações não detectadas podem violá-las.

Para cada ponderação, a família contém dois contrastes por cinco folds, totalizando $M=10$ limites marginais. Com $\gamma=0,05$, utiliza-se
\begin{equation}\label{eq:bonferroni}
 [L_k,U_k]=\widehat\delta_k\ \pm\
 z_{1-\gamma/(2M)}\widehat{\operatorname{SE}}_k,
\end{equation}
em que $z_q$ é o quantil $q$ da normal padrão. Os limites combinados são
\begin{equation}\label{eq:combinacaolimites}
 L=\sum_k\omega_kL_k,\qquad U=\sum_k\omega_kU_k,
 \qquad\omega_k=\frac{\sum_{g\in k}N_g}{\sum_gN_g}.
\end{equation}
A desigualdade da união de Bonferroni dispensa independência entre folds, mas a cobertura nominal depende dos limites marginais e do modelo auxiliar. Cada análise condiciona separadamente no seu treino e par de modelos, pois os demais treinos contêm faixas que são teste em outro fold.

O alvo combinado é a média de riscos condicionais dos pares ajustados, sem repetir treino, seleção ou construção dos dados. A Tabela~\ref{tab:incerteza} não cobre a variabilidade do algoritmo integral, nem fundamenta inferência de catálogo ou tratamento dos folds como réplicas independentes \cite{bengio2004,bates2024}.

\end{document}
